\documentclass[11pt]{article}
\usepackage[a4paper,margin=27mm]{geometry}
\usepackage{amsmath,amssymb,amsthm,mathrsfs,hyperref}
\setlength{\emergencystretch}{2em}
\newtheorem{theorem}{Theorem}
\newtheorem{lemma}[theorem]{Lemma}
\newtheorem{proposition}[theorem]{Proposition}
\newcommand{\MA}{\operatorname{MA}}
\newcommand{\EE}{\mathcal E^1_\theta}
\title{Absolute finite-model estimates and genuine variational attainment\protect\\
under the explicit outer-output foundations}
\author{Unreviewed conditional research, continuation 200}
\date{6 October 2026}
\begin{document}
\maketitle
\begin{abstract}
On actual finite nef models, a vertical intersection graph proves
that taking the absolute value contracts every mixed Hodge seminorm.
This gives an absolute integral estimate, together with a quantitative
midpoint energy gap. Under the explicitly unproved genuine scalar,
moving-system and ordered-tail foundations, both estimates transfer
to genuine finite-energy functions. Arbitrary-target output
approximation then yields strong absolute integration against every
fixed intrinsic finite-dual-energy target. Model near maximizers are
strongly Cauchy; finite divisorial supporting planes exclude an energy
jump at a genuine weak cluster. The target objective is attained,
and both stationarity signs give its equation. Entire supporting
planes and divisorial target probes separate every pair of solutions.
These conclusions remain conditional on all stated foundations.
Identification of all unbounded mixed products, full negligible-set
nullness and strict Borel comparison are not proved here, so original
P16.5 remains open.
\end{abstract}

\section{Finite geometry: an absolute value contraction}
Let $A$ be an actual nef rational bundle on a smooth projective
$n$-fold, $n\geq1$, over an algebraically closed trivially valued
characteristic-zero field. Work first only with bounded rational
relatively nef model scalars on its original analytic carrier.
On a smooth common determination $\mathcal X$ the central fiber is
$F=\sum_i b_iE_i$ with SNC support; all $b_i$ are retained.
Write $T_\xi=(A,\xi)$ for a nef model slot,
$e_A(\xi)=T_\xi^{n+1}/(n+1)$ and
$J_A(\xi)=V_A\xi(v_0)-e_A(\xi)$, $V_A=A^n$.
The raw measure has coefficients $b_i(T_\xi|_{E_i})^n$ and mass
$V_A$. For a vertical scalar $f$ define
\[
 Q_\Gamma(f)=-(0,f)^2\Gamma,\qquad
 \Gamma=T_1\cdots T_{n-1},
\]
where the slots are nef on all central components.
For $n=1$ the product is empty. These are finite intersection
numbers; no unbounded ample estimate is imported.

\begin{lemma}[Graph contraction]\label{graph}
For a rational model function $f$, on a determination also realizing
the model function $|f|=\max(f,-f)$,
\[
 Q_\Gamma(|f|)\leq Q_\Gamma(f).                            \tag{1}
\]
More generally any scalar $1$-Lipschitz map on the finitely many
vertical values contracts the corresponding graph quadratic form.
\end{lemma}
\begin{proof}
If $f_i$ is the value at $v_{E_i}$, its vertical divisor is
$D_f=\sum_i b_if_iE_i$. Set $c_{ij}=E_iE_j\Gamma$ for $i\ne j$.
SNC effectivity and nef restrictions give $c_{ij}\geq0$.
The restriction of the full fiber $F$ to each central component is
numerically trivial, so $\sum_j b_jE_iE_j\Gamma=0$.
Expanding the diagonal coefficients with this relation gives
\[
 -D_f^2\Gamma=\sum_{i<j}b_ib_jc_{ij}(f_i-f_j)^2.            \tag{2}
\]
Pulling slots and divisors to a further determination preserves
their finite intersections. The absolute-value model exists on an
actual vertical-ideal blowup: clear the denominators of $D_f$,
multiply the two fractional invertible ideals representing $f,-f$
by one sufficiently large power of the base parameter, take their
sum, and principalize it. The scalar logarithm of the norm of their
ideal sum is the maximum of the two scalar logarithms. Undoing the common parameter power
gives $\max(f,-f)$, with its value $|f_i|$ at every new divisorial
component. Characteristic-zero SNC resolution is ordinary algebraic
background; the generic fiber and original model scalars are retained.
Apply (2) on this common refinement and use
$||f_i|-|f_j||\leq|f_i-f_j|$. This proves (1).
The last assertion concerns the finite form (2); it asserts model
realizability only when the transformed values come from a model.
\end{proof}

Polarization of (2) gives Cauchy--Schwarz for $Q_\Gamma$.
The correct tangent identity, with $S=(0,\xi-\eta)$, is
\[
 T_\eta^{n+1}-T_\xi^{n+1}+(n+1)ST_\eta^n
 =-\sum_{j=0}^{n-1}(n-j)S^2T_\xi^jT_\eta^{n-1-j}.          \tag{3}
\]
Indeed telescope the first difference, subtract its tangent term,
and telescope each remaining difference once more; the term indexed
by $j$ occurs for $k=j+1,\ldots,n$.
In particular
\[
 J_A(\xi)=\frac1{n+1}\sum_{j<n}(n-j)
                     Q_{T_\xi^jT_0^{n-1-j}}(\xi).         \tag{4}
\]
The historically printed coefficient $j+1$ is not used with these
factor positions. All coefficients in (4) are positive.

\section{Finite slot estimates and an absolute integral}
Put $d_A(\xi,\eta)=\max_{j<n}Q_{T_\xi^jT_\eta^{n-1-j}}(\xi-\eta)$.
For completeness the finite slot-change argument is recalled.
Concavity of the finite energy follows by differentiating its
polynomial on a nef segment: its second derivative is
$-nQ_{T_{\xi_t}^{n-1}}(\xi-\eta)\leq0$.
The tangent remainder is therefore convex in its first variable;
by (3) it is comparable, with dimension constants, to $d_A$.
Expanding the average of two slots gives
\[
 d_A(\xi,\xi')\asymp_n
 Q_{T_{(\xi+\xi')/2}^{n-1}}(\xi-\xi').                    \tag{5}
\]
For $f=\xi-\xi'$, $\tau=(\xi+\xi')/2$, set
$a=d_A(\xi,\xi')$, $b=\max(d_A(\xi,\eta),d_A(\xi',\eta))$,
and $b_j=Q_{T_\eta^jT_\tau^{n-1-j}}(f)$.
Then $b_0\leq C_na$ and $b_{n-1}\leq4b$.
In $b_{j+1}-b_j$ split one occurrence of $f$ into
$T_\xi-T_\xi'$ and apply Cauchy--Schwarz separately with these
two nef slots. Since $2T_\tau-T_\xi=T_\xi'$ is nef and conversely,
their $f$ seminorms are at most $2b_j$; their other seminorms
are at most twice $d_A(\tau,\eta)\leq C_nb$ by convexity of the
tangent remainder. Thus
\[
 |b_{j+1}-b_j|\leq C_n\sqrt{b_jb}.
\]
Induction for $a\leq b$ gives
$b_j\leq C_na^{2^{-j}}b^{1-2^{-j}}$; for $a>b$ the final
bound $b_{n-1}\leq4b$ gives the same desired last-slot estimate.

To obtain the quasi-triangle estimate, split $\xi-\eta$ through a
third potential $\zeta$ in the average seminorm (5). Apply the
last-slot bound to the two differences. Convexity bounds the distances
from the averaged slot to $\zeta$ by
$C_n\max(d_A(\xi,\zeta),d_A(\eta,\zeta))$; its distances to
$\xi,\eta$ are bounded by $C_nd_A(\xi,\eta)$ by expansion.
Writing $d=d_A(\xi,\eta)$ and
$b=\max(d_A(\xi,\zeta),d_A(\eta,\zeta))$, this gives
$d\leq C_nb^{2^{1-n}}\max(d,b)^{1-2^{1-n}}$.
Divide by the power of $d$ when $d\geq b$; otherwise the conclusion
is immediate. Hence $d\leq C_nb$.
Average an arbitrary finite list of nef slots and expand its
nonnegative multinomial terms. The preceding last-slot bound,
quasi-triangle estimate and (4) then give
\[
 Q_{T_\psi^jT_0^{n-1-j}}(\xi-\eta)
 \leq C_n d_A(\xi,\eta)^{2^{1-n}}
       M^{1-2^{1-n}},\quad
 M=\max(J_A(\xi),J_A(\eta),J_A(\psi)).                    \tag{6}
\]
For $n=1$ there is one empty-slot seminorm and the same result
follows directly. Constants in the scalars affect neither these
seminorms nor $J_A$.

\begin{theorem}[Absolute raw finite estimate]\label{finiteabs}
For bounded rational relatively nef model scalars $\xi,\eta,\psi$,
with $\alpha=2^{-n}$,
\[
 \begin{split}
 \int|\xi-\eta|\,d\MA_{A,\mathrm{raw}}(\psi)
 \leq{}&V_A|\xi(v_0)-\eta(v_0)|\\
 &+C_n d_A(\xi,\eta)^\alpha J_A(\psi)^{1/2}
              M^{1/2-\alpha}.
 \end{split}                                             \tag{7}
\]
This includes $V_A=0$. There is no volume denominator.
\end{theorem}
\begin{proof}
Use the actual absolute-value model from Lemma~\ref{graph} and
telescope $T_\psi^n-T_0^n$. The difference between the integral
on the left and its trivial-model integral is
\[
 \sum_{j=0}^{n-1}(0,|\xi-\eta|)(0,\psi)T_\psi^jT_0^{n-1-j}.
\]
Cauchy--Schwarz, (1), (6), and (4) bound its absolute value by
the second line of (7). The trivial measure is $V_A\delta_{v_0}$.
If an energy factor vanishes, use the seminorm formulas directly
or a finite polynomial limit. In dimension one the last exponent
is zero and the bound is $C_1d_A^{1/2}J_A(\psi)^{1/2}$.
If $V_A=0$, positivity and raw mass force the measure to be zero.
\end{proof}

\section{A quantitative finite midpoint defect}
Define the raw finite pairing
\[
 I_A(\xi,\eta)=\int(\xi-\eta)\,
       d(\MA_A(\eta)-\MA_A(\xi))
   =\sum_{j<n}Q_{T_\xi^jT_\eta^{n-1-j}}(\xi-\eta)\geq0.
\]
Thus $d_A\leq I_A$. Let $g(t)=e_A((1-t)\eta+t\xi)$.
Its second derivative is the negative of $n$ times
\[
 \sum_{j<n}\binom{n-1}{j}t^j(1-t)^{n-1-j}
                       Q_{T_\xi^jT_\eta^{n-1-j}}(\xi-\eta).
\]
The midpoint defect equals $\int_0^1G(t)(-g''(t))\,dt$, where
$G(t)=t/2$ for $t\leq1/2$ and $G(t)=(1-t)/2$ otherwise.
On $[1/4,3/4]$, $G\geq1/8$ and each power product is at least
$4^{-(n-1)}$. Consequently
\[
 e_A((\xi+\eta)/2)-\frac{e_A(\xi)+e_A(\eta)}2
       \geq c_n I_A(\xi,\eta),\qquad c_n=n4^{-(n+1)}>0.    \tag{8}
\]
This is elementary finite polynomial calculus, including nonreduced
fibers through (2). The constant need not be optimal.

\section{Explicit genuine antecedents; none is an accepted grant}
Now let $L$ be any actual big rational bundle on smooth projective $X$
over any retained field. Put $\theta=c_1(L)$, $K=X^{\mathrm{an}}$ and
$V=\operatorname{vol}(L)>0$. The genuine space uses one all-point
decreasing FS net on varying actual rational bundles, with classes
tending to $\theta$. All remaining assertions are conditional on
the following proposed foundations. Supported semivaluations remain.
\begin{enumerate}
\item Genuine closure, convexity, envelopes and scalar section energy
are well defined. For $R=P_\theta(0)$, normalized $u\in\EE$ satisfy
$u\leq R$, $E(R)=0$ and $B(u)=-E(u)\geq0$. Energy is concave,
decreasing continuous and weakly usc, with $E(u+c)=E(u)+cV$.
The actual extension is $E(u)=\inf_{f\in C(K),\ f\geq u}\mathscr E(f)$,
with $E(P_\theta(f))=\mathscr E(f)$ for continuous obstacles.
The normalized genuine cone is compact in the weak topology induced
by its finite divisorial values; these coordinates determine all-point
order. Convex combination with a fixed function is weakly continuous.
Every such $u$ has nonpositive rational model majorants
$f_i\downarrow u$, $f_i(v_0)=0$, whose envelopes $u_i=P_\theta(f_i)$
decrease with energies to $u$. Perturbed envelopes decrease at each
fixed original continuous direction. For one actual Kodaira divisor
$Z$, $R$ and all bounded-obstacle envelopes are finite continuous on
$U=K\setminus Z$. A genuine finite-energy pole on $Z$ exists.
\item Simultaneous complete full-system resolutions for any fixed
finite list of model obstacles share an actual factor
$p_m^*L=A_m+T_m$, with $T_m$ effective and $mA_m$ generated.
Their section stages have $\phi_{m,f}=q_{T_m}+\psi_{m,f}$ and
$\phi_{m,0}=q_{T_m}$. The $\psi_{m,f}$ are rational relatively
generated model scalars on the same nef $A_m$. We have $A_m^n\to V$,
$e_{A_m}(\psi_{m,f})\to E(P_\theta(f))$, generic values tending to
zero, and pushed raw measures converging to the finite divisorial model
output. Bounded quotients converge uniformly on compact subsets of $U$
to $P_\theta(f)-R$. Divisorial lifts are unique. The inequalities
$\phi_{m,f}\leq P_\theta(f)$ hold on the entire original carrier.
Auxiliary energies equal the full-section traces of their relatively
generated nef systems.
\item The target-independent map $M$ has positive raw Radon mass $V$,
constant invariance and algebraic nullness. At a model envelope it is
the finite divisorial derivative of $\mathscr E(f)=E(P_\theta(f))$.
Canonical majorants have $M(u_i)\to M(u)$ weakly. Every actual test
has finite relative integrals at outputs. For each fixed normalized
$v\in\EE$ and cost bound $B$, its majorants satisfy
\[
 \sup_{B(w)\leq B}\int(v_i-v)\,dM(w)\longrightarrow0.       \tag{9}
\]
Also there are nonnegative $a_\epsilon\in C(K)$, zero on $Z$ and
below $b_v=R-v$ off $Z$, with
\[
 \sup_{B(w)\leq B}\int|b_v-a_\epsilon|\,dM(w)<\epsilon.     \tag{10}
\]
These are the conditional ordered-gap and fixed-collar conclusions of
the completed partial paper. They do not include strong integration
against an arbitrary prescribed target.
\item Quadratic genuine contraction holds:
$B((1-r)R+ru)\leq C_nr^2B(u)$, $0<r<1$.
For a fixed rational model direction $h$ and normalized nonpositive
model obstacles of cost at most $B$, the original increment estimate
holds uniformly over obstacles and refinements:
\[
 \left|\mathscr E(f+th)-\mathscr E(f)-t\int h\,dM(P_\theta(f))\right|
 \leq K_ct^2/2+|t|A_h(B)/c,                               \tag{11}
\]
for fixed $c>B/V$ and $|t|<t_c$. Here $K_c,t_c$ are independent
of $f$, and $A_h(B)=\|h\|_\infty(2B+2V\|h\|_\infty)$ suffices.
Original model functions are uniformly dense in $C(K)$.
\end{enumerate}
These are unproved antecedents from the scalar/outer-output and
ordered-tail drafts. They assume no genuine strong completeness,
arbitrary-target strong integration, attainment or uniqueness.
Physical polar section stages are not declared original $\EE$.
The finite model conclusions of Sections1--3 are separate.

\section{Absolute estimates and midpoint transfer on the genuine space}
For normalized genuine endpoints define
\[
 I(u,v)=\int(u-v)\,d(M(v)-M(u)).                           \tag{12}
\]
Relative integrability makes it finite, without absolute reference
integrability. First let $a,b,w$ be normalized model envelopes.
Under item2, the bounded differences $\psi_{m,a}-\psi_{m,b}$ converge
uniformly on compact subsets of $U$ to $a-b$. Fixed continuous collars
and algebraic nullness pass their absolute integrals to the finite
divisorial limiting output. The same argument at each of the two
endpoint outputs passes the four pairings defining $I_{A_m}$, so
$I_{A_m}(\psi_{m,a},\psi_{m,b})\to I(a,b)\geq0$.
Generic errors in (7) vanish only after the moving-degree limit;
the $J$ energies converge to $B(a),B(b),B(w)$. Since $d_{A_m}\leq I_{A_m}$,
\[
 \int|a-b|\,dM(w)\leq C_n I(a,b)^\alpha B(w)^{1/2}
                  \max(B(a),B(b),B(w))^{1/2-\alpha}.       \tag{13}
\]
For arbitrary genuine $u,v,w$, take canonical majorants
$a_i\downarrow u$, $b_j\downarrow v$, $w_l\downarrow w$.
Equation(9) controls each changing gap uniformly over these outputs.
For fixed tests, (10) passes their relative integrals weakly.
All four pairings in $I(a_i,b_j)$ therefore converge jointly to (12).
In particular $I(u,v)\geq0$. On $U$,
\[
 ||a_i-b_j|-|u-v||\leq(a_i-u)+(b_j-v).
\]
Its right side has uniformly vanishing output integrals by (9).
The fixed $|u-v|=|b_u-b_v|$ has the continuous approximant
$|a_{\epsilon,u}-a_{\epsilon,v}|$ from (10), with the sum of the two
absolute errors. Its fixed weak integral consequently passes too.
Taking the three indices in (13) proves
\[
 \int|u-v|\,dM(w)\leq C_n I(u,v)^\alpha B(w)^{1/2}
                  \max(B(u),B(v),B(w))^{1/2-\alpha}.       \tag{14}
\]
This uses the stated ordered-tail package, not a transferred ample
unbounded estimate.

The midpoint requires an actual trace inequality. For fixed model
envelopes $a,b$ put $z=(a+b)/2$ and
$\rho_m=(\psi_{m,a}+\psi_{m,b})/2$ on $A_m$.
Then $q_{T_m}+\rho_m\leq z$. Fix any original $f\in C(K)$, $f\geq z$.
In every cleared degree $\ell$, multiply the entire auxiliary section
space by the effective section of $T_m$. Normal proper birational
descent injects it into the original $H^0(X,\ell L)$.
Every auxiliary section weighted below $\rho_m$ lies below $f$ after
this multiplication, since its scalar lies below
$q_{T_m}+\rho_m\leq f$ on the whole supported carrier.
Complete this injected adapted flag to a full original flag.
Complementary original weights are at least $\ell\min_Kf$; the
canonical norms of regular sections are at most one. Integer flooring
causes only the vanishing degree error.

At fixed $m$ the auxiliary full trace is $e_{A_m}(\rho_m)$:
the rational tensor average is relatively generated; twisting by a
high base degree makes it generated on the compactification.
Ordinary generated-nef volume/intersection equality and the lattice
intercept give this trace identity. Thus first $\ell\to\infty$ gives
\[
 \mathscr E(f)\geq e_{A_m}(\rho_m)
                           +(\min_K f)(V-A_m^n).          \tag{15}
\]
The continuous majorant and its minimum are fixed before $m$.
Send $m\to\infty$, then take the outer infimum over $f\geq z$.
It gives $E(z)\geq\limsup_me_{A_m}(\rho_m)$; the lost-volume term
has been retained until its limit. Equation(8), the endpoint energy
limits and the pairing limits imply
\[
 E((a+b)/2)-\frac{E(a)+E(b)}2\geq c_nI(a,b).               \tag{16}
\]
Take decreasing canonical majorants of arbitrary $u,v$.
Genuine convexity and decreasing energy continuity pass all three
energies, and the pairing convergence above passes $I$.
Thus (16) holds for all normalized genuine $u,v\in\EE$.
This argument never assigns original finite energy to a physical
polar complete-system stage.

\section{Entire supporting planes}
For a model obstacle $f$, concavity and its finite derivative give
$E(v)\leq\mathscr E(f)+\int(v-f)d\mu_f$ for every genuine $v$.
First apply the continuous plane to continuous majorants of $v$, then
pass their finite divisorial values. Those values are finite and
converge. Apply the result to $a=P_\theta(f)$, with
$E(a)=\mathscr E(f)$ and $a\leq f$. It forces
$\int(f-a)d\mu_f=0$, so the whole plane at the model envelope is
$E(v)\leq E(a)+\int(v-a)dM(a)$.
For canonical $a_i\downarrow u$, (9) gives
$\int(a_i-u)dM(a_i)\to0$. The fixed $v-u$ integral passes to $M(u)$
by (10), separately for its two relative costs. Hence
\[
 E(v)\leq E(u)+\int(v-u)\,dM(u),\qquad u,v\in\EE.         \tag{17}
\]
Restore constants by the raw mass law. Using (14) at either endpoint
of (17) also gives, for normalized $u,v$,
\[
 |E(u)-E(v)|\leq C_n I(u,v)^\alpha
                        \max(B(u),B(v))^{1-\alpha}.       \tag{18}
\]

\section{Every fixed intrinsic target has strong absolute integration}
Let $\nu$ have raw mass $V$, give reference poles zero mass, integrate
all genuine normalized relative costs and have finite dual energy
\[
 D_\nu=\sup_{u\in\EE,\ u(v_0)=0}
                  \{E(u)+\int(R-u)d\nu\}<\infty.          \tag{19}
\]
The genuine finite-energy pole on $Z$ makes $\nu(Z)=0$.
If the intended intrinsic class is defined just by all-test
integrability, its equivalence with (19) is a separate scalar/coercivity
foundation; it is not silently assumed in this paper.

We recall, and justify, the target approximation needed here.
For canonical model majorants of any fixed test, their costs
$R-u_i$ increase continuously on $U$. Radon regularity gives
$\int\sup_i(R-u_i)d\nu=\sup_i\int(R-u_i)d\nu$ even for directed
nets: each compact subset of an open positive superlevel has a finite
cover by stage superlevels, and directedness chooses one dominating
index. Finite lower layer-cake sums and truncation give the integral
identity. This is not measurable monotone convergence for arbitrary nets.
Together with decreasing energy continuity, model objectives approach
every genuine objective. Choose normalized model envelopes $s_k$
with objective at least $D_\nu-1/k$.

Quadratic contraction gives
$I_\nu(u)\leq B(u)/2+2C_nD_\nu$, by testing the dual supremum at
$(1-r)R+ru$, $r=1/(2C_n)$.
Consequently $B(s_k)\leq B_\nu:=4C_nD_\nu+2$.
Estimate(11) passes to every genuine base with the same constants:
use canonical model majorants and pass the entire increment at fixed
$t$, their decreasing perturbed energies and their weak output measures.
For $z=P_\theta(s_k+th)$, the obstacle inequality gives
$I_\nu(z)\geq I_\nu(s_k)-t\int h\,d\nu$.
Near maximality yields
$E(z)-E(s_k)-t\int h\,d\nu\leq1/k$.
The uniform genuine increment estimate for both signs therefore gives
\[
 \left|\int h\,d(M(s_k)-\nu)\right|
 \leq\frac1{k|t|}+K_c|t|/2+A_h(B_\nu)/c.                 \tag{20}
\]
First $k\to\infty$, then $t\to0$ at fixed $c$, then $c\to\infty$.
Uniform model density and raw mass give $M(s_k)\to\nu$ weakly.
The perturbed obstacle is the actual $s_k+th$, not a larger physical
model majorant. No attainment was used in this approximation.

For each fixed $v$, extend $b_v=R-v$ by zero on $Z$. It is globally
nonnegative lsc: $R$ is continuous on $U$, $v$ is usc, and the assigned
value on $Z$ is zero. Thus $b_v-a_\epsilon$ from (10) is nonnegative
lsc. Portmanteau and (20) give
$\int(b_v-a_\epsilon)d\nu\leq\epsilon$ as well.
Since $\nu(Z)=0$, this is the actual relative integral.
The fixed $|u-v|=|b_u-b_v|$ has the original continuous approximant
$|a_{\epsilon,u}-a_{\epsilon,v}|$, with error at most the two
cost errors, both at the output approximants and at $\nu$.
Therefore its integral passes weakly to $\nu$.
Apply (14) with output $M(s_k)$ and $B(s_k)\leq B_\nu$ to obtain,
for every normalized genuine pair $u,v$,
\[
 \int|u-v|\,d\nu\leq C_n I(u,v)^\alpha B_\nu^{1/2}
                    \max(B(u),B(v),B_\nu)^{1/2-\alpha}.   \tag{21}
\]
The target is fixed, and this estimate holds for every pair with the
same $B_\nu$. Hence it gives strong absolute integration for changing
pairs whose costs are bounded and whose $I$ tends to zero. This is
stronger than the earlier fixed-test output limit alone.

\section{A genuine maximizing potential, with no energy jump}
Use the same normalized model near maximizers $s_k$.
Relative integrals are affine on genuine convex combinations. Thus
the midpoint objective is its endpoint objective average plus the
midpoint energy defect. Its value is at most $D_\nu$, so (16) gives
\[
 c_nI(s_k,s_l)\leq\frac1{2k}+\frac1{2l}.                  \tag{22}
\]
Their costs are bounded by $B_\nu$, and (18) makes $E(s_k)$ Cauchy.
Write $E_*\geq-B_\nu$ for its limit. Compactness in item1 supplies
a genuine normalized weak cluster $u$. Weak usc gives
$E(u)\geq E_*$, so $u\in\EE$ with $B(u)\leq B_\nu$.

We need equality, not just this inequality. For each fixed $j$,
$M(s_j)$ has finite divisorial support. Its values are finite
continuous coordinates of the weak topology. Along the cluster subnet,
\[
 \int(s_k-s_j)dM(s_j)\longrightarrow\int(u-s_j)dM(s_j).
\]
The absolute bound(14) and (22) control the integral uniformly for
all sufficiently late $k$, by a quantity tending to zero as $j\to\infty$.
Apply the entire supporting plane(17) at the fixed model $s_j$:
\[
 E(u)\leq E(s_j)+\int(u-s_j)dM(s_j)\leq E(s_j)+o_j(1).
\]
Let $j\to\infty$. This proves $E(u)\leq E_*$ and hence equality.
The finite support, the fixed plane and the order of limits exclude
the energy jump without strong completeness being assumed.

Along the same subnet, $(s_k+u)/2\to u$ weakly. Weak usc bounds
its limiting energy above by $E(u)$; concavity bounds it below by
$(E(s_k)+E(u))/2\to E(u)$. Its midpoint defect therefore tends to zero.
Equation(16) gives $I(s_k,u)\to0$ along that subnet.
The fixed-target absolute estimate(21) now gives
$\int|s_k-u|d\nu\to0$. Hence the relative target integrals converge,
as do the energies. Their near-maximal objectives converge to the
objective of the actual genuine $u$, proving
\[
 E(u)+\int(R-u)d\nu=D_\nu.                               \tag{23}
\]
This is attainment at an original genuine function, not an abstract
completion point or a compact cluster with unchecked target integral.

For this maximizer the envelope inequality used in (20), with no
near-maximizer error, gives
$E(P_\theta(u+th))-E(u)-t\int h\,d\nu\leq0$.
Both genuine derivative signs obtained from (11) imply
$\int h\,dM(u)=\int h\,d\nu$ for each rational model $h$.
Density and equal raw mass give
\[
                              M(u)=\nu.                  \tag{24}
\]
All statements in this section remain conditional on items1--4.

\section{Separation of every solution by actual divisorial probes}
First, each output lies in the target class(19). Its raw mass is $V$,
it excludes reference poles by algebraic nullness, and it integrates
every relative test by item3. The whole plane(17) shows
\[
 E(v)+I_{M(u)}(v)\leq E(u)+I_{M(u)}(u)\quad(v\in\EE).
\]
Thus its dual energy is finite and equals the finite expression on
the right. Conversely (24) provides a solution for every target(19).

Let normalized genuine $u,v$ have the same output $\nu$.
By the whole planes, both maximize the same objective. Its midpoint
cannot exceed the supremum; its affine integral and concave energy
make the midpoint energy defect zero. Hence (16) gives $I(u,v)=0$.
This alone has not yet asserted pointwise equality.

Fix any rational divisorial valuation $z$. Its coordinate is finite
continuous on the compact normalized genuine cone, so all normalized
genuine values there have one uniform finite lower bound.
The reference value $R(z)$ is finite. Consequently the actual raw
target $\nu_z=V\delta_z$ excludes reference poles, integrates every
relative test, and has
\[
 0\leq D_{\nu_z}\leq V\sup_{a(v_0)=0}(R(z)-a(z))<\infty.
\]
It is an intrinsic target, rather than a formal abstract probe.
By (24) there is an actual genuine $w_z\in\EE$ with
$M(w_z)=V\delta_z$. Its energy cost is finite.
Use (14) at this actual output and $I(u,v)=0$ to obtain
\[
 V|u(z)-v(z)|=\int|u-v|dM(w_z)=0.
\]
Every rational divisorial coordinate agrees. All-point order
determination in item1 gives $u=v$ on the entire original carrier,
including supported semivaluations. Removing normalization gives
uniqueness modulo constants for \emph{every} solution, not just the
one maximizer constructed by a selected near-maximizer sequence.

\begin{theorem}[Conditional variational bijection]
If all the explicitly stated items1--4 hold in their entire genuine
scope, the proposed outer-output map gives a bijection from
$\EE/\mathbb R$ onto the intrinsic raw mass-$V$ targets(19).
Its solutions are attained genuine potentials, unique modulo constants.
The assertion uses no assumption of genuine strong completeness or
pre-existing target attainment.
\end{theorem}
\begin{proof}
The image inclusion follows from the whole planes and integrability.
Existence is (23)--(24). The divisorial-probe argument just given
proves injectivity for every solution. Constant invariance is item3.
\end{proof}

\section{What this conditional result does not establish}
The antecedents are substantial mathematical obligations.
Genuine compactness and scalar energy, the entire target-independent
output construction, common actual generated moving systems,
all ordered-tail estimates and the uniform original increment estimate
still require independent verification. Historical source labels,
compiled papers or another unreviewed manuscript do not establish them.
The finite graph and polynomial statements themselves are limited to
actual bounded rational nef models.

Even if those antecedents and the resulting variational bijection
were verified, additional parts of original P16.5 remain:
\begin{itemize}
\item Identify the proposed $M$ with the genuine full mixed
Monge--Amp\`ere operator and every admissible simultaneous unbounded
moving-product and relative weighted limit.
\item Prove every required nonalgebraic negligible-set nullness
statement, without replacing it by algebraic-divisor nullness.
\item Prove locality and comparison on the original strict Borel
winning sets. For usc $u,v$, $\{u<v\}$ need not be open; no weak
Portmanteau argument above passes its indicator.
\item Verify the intrinsic class equality if the original class is
specified by all-test integrability alone, and retain the exact
additional zero-volume range from the original task. This paper only
treats the positive-volume domain.
\end{itemize}
In particular, the original model envelope of a maximum of larger
physical obstacles need not be the maximum of the two genuine
envelopes. Such an identity has not been used as a locality proof.
The actual trace injection(15) is only an energy lower bound;
it does not identify all products or erase a polar common factor.

This manuscript is unreviewed conditional progress, not a full solution
of original P16.5. It changes neither the solved-original count nor
the independently accepted supporting-paper count. The active native
review and author campaigns received no hot feedback; their private
outputs were not read. The user's current editor remains unchanged.

\begin{thebibliography}{9}\small
\bibitem{BJ} S. Boucksom and M. Jonsson,
\emph{Global pluripotential theory over a trivially valued field},
\href{https://arxiv.org/pdf/1801.08229v4}{arXiv:1801.08229v4},
finite-model intersection and Hodge calculus. The orientation of the
tangent coefficients is corrected directly in (3).
\bibitem{finite} Whole 10-page unreviewed corrected finite-model and
Radon-squeeze paper, 6 October 2026. The finite slot argument is
reconstructed in Section2 here rather than supplied as an accepted grant.
\bibitem{foundation} Whole 14-page unreviewed genuine scalar,
section-energy and outer-output foundation paper, 6 October 2026.
\bibitem{tail} Whole completed 12-page unreviewed
\emph{Bounded-obstacle section approximation and ordered energy-gap tails},
6 October 2026; its package F is explicitly unaudited.
\bibitem{density} Conditional bounded-obstacle full-section density
and general bounded weighted-limit draft, continuation198.
\bibitem{target} Conditional arbitrary-target output approximation
before attainment, continuation199, 6 October 2026.
\end{thebibliography}
\end{document}
